\documentclass[a4paper,12pt,oneside]{article}
\makeatletter
\def\input@path{{../../}}
\makeatother
\usepackage{discrete-practice}

\begin{document}

\PracticeHeader{2}{4}{Деревья. Расстояния}

\Task{Задание 1}
Найти радиус, диаметр, центр(ы) графов ниже.
\begin{numlist}
  \item \raisebox{-.5\height}{\includegraphics[width=.74\linewidth,height=.14\textheight,keepaspectratio]{Practice_4_1.png}}
  \item \raisebox{-.5\height}{\includegraphics[width=.74\linewidth,height=.14\textheight,keepaspectratio]{Practice_4_2.png}}
  \item \raisebox{-.5\height}{\includegraphics[width=.74\linewidth,height=.14\textheight,keepaspectratio]{Practice_4_3.png}}
  \item \raisebox{-.5\height}{\includegraphics[width=.74\linewidth,height=.14\textheight,keepaspectratio]{Practice_4_4.png}}
  \item \raisebox{-.5\height}{\includegraphics[width=.74\linewidth,height=.14\textheight,keepaspectratio]{Practice_4_5.png}}
\end{numlist}

\clearpage

\Task{Задание 2}
\begin{numlist}
  \item Построить дерево имеющие 9 вершин, 3 центра, $\mathrm{rad}=3$,
    $\mathrm{diam}=6$. Продемонстрировать на них работу алгоритма.
  \item Постройте леса деревьев, где количество вершин равно: а) 4, б) 5, в) 6.
\end{numlist}

\Task{Задание 3}
\begin{numlist}
  \item Найти и выделить на графе все (обоснуйте свои ответы)
    \begin{itemize}
      \item точки сочленения
      \item блоки
      \item мосты
      \item компоненты реберной двусвязности (укажите размер максимальной и минимальной компоненты)
      \item компоненты вершинной двусвязности (укажите размер максимальной и минимальной компоненты)
      \item проверьте граф на двудольность
      \item найти радиус/диаметр/центр
      \item найти цикломатическое число и ФСЦ
      \item закодируйте полученное дерево
    \end{itemize}
  \item Постройте дерево блоков и точек сочленения.
  \item Постройте матрицу инцидентности.
\end{numlist}
\begin{center}
  \includegraphics[width=.56\linewidth]{Practice_4_6.png}
\end{center}

\end{document}
